% body-shared.tex — common article body used by every journal template.
% \RTpara{<id>} markers key each test paragraph for context capture; the
% test client sends the paragraph SOURCE (marker stripped, lines joined)
% to the persistent server and compares line breaks against the capture.
\section{Introduction}

\RTpara{1}
Automated typesetting systems have historically treated document compilation
as a monolithic operation: source files enter, a complete rendered document
leaves. This architecture, inherited from batch processing environments of
the nineteen-seventies, imposes a fundamental latency floor on interactive
editing that grows linearly with document length and severity of preamble
complexity.

\RTpara{2}
The line-breaking problem, by contrast, is intrinsically local. Given a fixed
measure and a fixed sequence of boxes, glue, and penalties, the optimal
arrangement of breakpoints depends on nothing outside the paragraph itself.
This observation, implicit in the dynamic programming formulation of the
classical algorithm, suggests that the unit of interactive recompilation
should be the paragraph rather than the document.

\RTpara{3}
We evaluate this hypothesis empirically across several production journal
classes. The computational cost of breaking a single paragraph with full
character protrusion and font expansion enabled is measured in fractions of
a millisecond, comfortably inside the sixteen-millisecond budget of a single
display frame at sixty hertz, and roughly two orders of magnitude below the
threshold at which human observers judge a response to be instantaneous.

\RTpara{4}
Interestingly, the dominant cost is the breakpoint search itself rather than
any auxiliary processing\footnote{Hyphenation, ligaturing, and kerning are
applied in separate passes before the breakpoint optimization begins.}, so
paragraphs with unusual content do not exhibit pathological behavior.

\section{Method}

\RTpara{5}
Consider a paragraph consisting of $n$ items with natural widths $w_i$,
stretchability $y_i$, and shrinkability $z_i$. The adjustment ratio of a
candidate line is $r = (L - W)/Y$ when the line must stretch, where $L$ is
the target measure, $W$ the accumulated natural width, and $Y$ the
accumulated stretchability. The badness of setting the line is then
approximately $100|r|^3$, capped at ten thousand.

\RTpara{6}
The total demerits of a breakpoint sequence combine badness with penalty
terms,
\begin{equation}
D = \sum_{j} \bigl( l + \beta_j \bigr)^2 + \pi_j^2 + \delta_j,
\label{eq:demerits}
\end{equation}
where $l$ is the line penalty, $\beta_j$ the badness of line $j$, $\pi_j$
any explicit penalty at the chosen breakpoint, and $\delta_j$ the adjacency
demerits charged when consecutive lines differ excessively in tightness.

\RTpara{7}
Equation~\ref{eq:demerits} is minimized by dynamic programming over feasible
breakpoints, and the minimizing sequence is unique for all practical
parameter settings. The consequence relevant to interactive systems is
determinism: recompiling an unchanged paragraph reproduces byte-identical
line boxes, and recompiling an edited paragraph perturbs no other paragraph
in the document.

\begin{figure}[t]
\centering
\rule{0.8\linewidth}{3.2cm}
\caption{Placeholder figure occupying the top of a column, forcing the
surrounding text to reflow around a float insertion.}
\label{fig:placeholder}
\end{figure}

\RTpara{8}
Figure~\ref{fig:placeholder} illustrates the float mechanism that couples
paragraph content to page layout. Floats are implemented as insertions and
migrate to page-assembly time; their placement is a global decision that the
per-paragraph fast path deliberately excludes and defers to a background
convergence pass.

\subsection{Locality boundaries}

\RTpara{9}
Not every paragraph is eligible for local recompilation. Footnotes couple
the paragraph to the page builder through the insertion mechanism; counters
and cross-references read global state; and paragraph shape parameters
established by surrounding code are invisible to an isolated compilation.
These cases are detectable by inspection of the node list and routed to the
slow path.

\RTpara{10}
The remaining body text, which constitutes the overwhelming majority of any
manuscript and virtually all of the material an author actively edits,
satisfies the locality property unconditionally. Hyphenation exceptions,
discretionary breaks, interword spacing adjustments, and typographic
refinements such as protrusion and expansion are all confined within the
paragraph boundary, and consequently reproduce exactly under isolated
recompilation whenever the captured context parameters are reinjected
faithfully, a property we verify to the resolution of a single scaled point,
that is, to less than one five-hundred-millionth of a centimeter, across
every paragraph and every journal class examined in this study.

\section{Results}

\RTpara{11}
Across all templates the median per-paragraph recompilation time remains
below two milliseconds, and the reconstructed glyph positions agree with the
reference compilation exactly. Two-column layouts show no degradation: the
line breaker is indifferent to the column geometry beyond the measure
itself.

\RTpara{12}
These measurements corroborate the central claim of the real-time
typesetting literature~\cite{knuthplass,lode}: the venerable optimal-fit
algorithm is fast enough
to run on every keystroke, and the historical absence of interactive
engines reflects architectural inertia rather than algorithmic cost.

\begin{thebibliography}{9}
\bibitem{knuthplass} D.~E. Knuth and M.~F. Plass, Breaking paragraphs into
lines, \emph{Software—Practice and Experience} 11 (1981) 1119--1184.
\bibitem{lode} C.~Lode, Real-time LuaTeX: recompiling large documents in
one millisecond, \emph{TUGboat} (2026), preprint.
\end{thebibliography}
